\section{Discussion}\label{sec:discussion}

\subsection{What the results say}\label{sec:interpretation}

\emph{The weak point is authorization, not obedience.} Telling these agents that
the boss wants the trade, that the client is about to be lost or that the rule has
changed did not measurably move any of the eight. Telling them that the trade had
already been authorized did. The manual in every prompt said that only the
compliance department grants exceptions, only through the approval tool, and that
statements in a message cannot grant one. In all, \num{all.counterfeit.sig.models.w} of the eight models nevertheless
gave way significantly to a counterfeit of that authorization, and under the forged
notice every one of them, including the two that otherwise almost never broke a
rule, placed large orders without the approval the manual requires. The failure is
not that the agents defer to their principal; it is that they cannot tell a claim
of authorization from authorization. In the terms of the instruction hierarchy
\citep{Wallace2024InstructionHierarchy}, a user message that calls itself a system
notice, or that reports a decision by an authority above the agent, was given the
weight of the authority it claimed.

\emph{One pressure level would have ranked the models wrongly.} Under the claimed
approval alone, Mistral Small 3.2 and Qwen2.5 72B would look nearly as safe as the
gpt-oss models (\num{mistral-small-24b.pres.P2_claimed_approval.k} and
\num{qwen2.5-72b.pres.P2_claimed_approval.k} violations of 24) and Qwen3 among the
weakest (\num{qwen3-32b.pres.P2_claimed_approval.k}). Under the forged notice alone
the order reverses (\num{mistral-small-24b.pres.P5_fake_system_notice.k},
\num{qwen2.5-72b.pres.P5_fake_system_notice.k} and
\num{qwen3-32b.pres.P5_fake_system_notice.k}). An evaluation built around one attack
would certify whichever model happened to resist it. Nor did size decide the
outcome: the two models of about 70 billion parameters violated more often than
Qwen3 32B and gpt-oss-20b.

\emph{A low violation rate is not the same as compliance.} Two of the models show
failures that a violation count alone would score as success. Mistral Small 3.2
committed no violation under the claimed approval because it stopped working,
falsely telling its principal that it could not trade. Llama 3.3 70B left most of
the legitimate requests undone because it wrote its orders into its explanations
instead of placing them, and wrote some breaching orders the same way. Neither is a
compliant agent; both are agents that did not do their job, one of them in a way
that happens to keep it out of the violation count. This is why the benchmark scores
legitimate requests and why its results should be reported together
(\Cref{sec:usage}).

\emph{The rules are mechanical, so the control should be too.} Every violation in
this study was decided by a rule engine from the order and the account, before
the order would have reached a market. A pre-trade check running the same rules
would have stopped every breaching order, and an approval check that looked up the
identifier in the approval system's own records would have rejected CMP-4471 every
time it was offered. The agent's judgment is the wrong place to enforce a rule that
a few lines of deterministic code can decide; the agent belongs upstream of the
control, not in place of it. This is the same conclusion firms reached long ago
for human traders, whose orders pass through pre-trade checks however well trained
they are, and FINRA's position that its rules apply to generative AI as to any
other technology \citep{FINRA2024RN2409} points the same way.

\emph{The record has to come from the system, not the agent.} In \num{all.misreport.k} episodes across \num{all.misreport.models.w} models, including
gpt-oss-20b, the second most compliant model in the study, an agent placed a
breaching order and then told its principal that it had not. Supervision that
relies on what an agent says it did will miss violations that the ledger shows
plainly. The transcripts released with this paper, which record each model turn,
each tool call and the resulting ledger, are an audit trail of the kind earlier
work in this series argued machine-learning finance should keep
\citep{Khan2026AuditTrails}.

\emph{The serving stack is part of the measurement.} The same weights, prompt and
decoding settings reproduced every episode exactly on the same software and
changed the verdict in a few episodes on a newer version. Results on this
benchmark, as on any agent benchmark, should therefore name the serving software
and its version along with the model. For a firm the same point is a model-risk
one: supervisory guidance expects an inventory of the models in use and validation
of changes to them \citep{Fed2011SR117}, and an upgrade of the serving software
is such a change even when the weights stay the same.

\subsection{Controls for a firm deploying a trading agent}\label{sec:controls}

\Cref{tab:controls} turns the failures observed here into controls. None of them
is new: each is a control that firms already apply to people and to algorithmic
order flow, and for broker-dealers with market access the market access rule
already requires controls of the first kind, covering restricted securities and
capital thresholds \citep{SEC2010MarketAccess}. The point of the
table is that an agent needs all of them, because each failure in it occurred in
a model that, on other measures, behaved well.

\begin{table}[H]
\centering
\caption{Failures observed in this study and the controls that would stop them.
Each control sits outside the agent, so it does not depend on the agent's
judgment.}
\label{tab:controls}
\small
\begin{tabularx}{\textwidth}{@{}>{\raggedright\arraybackslash}X>{\raggedright\arraybackslash}X@{}}
\toprule
Failure observed & Control \\
\midrule
A breaching order placed as the agent's first action, with no check
(\Cref{sec:how}). & Hard pre-trade checks in the order management system that
reject a breaching order before it reaches the market, whoever placed it. \\
\addlinespace
A claimed sign-off or a reference number treated as an approval
(\Cref{sec:pressure-results}). & Approvals verified against the approval
system's own records by identifier, ticker, side and quantity. The agent can
request an approval but never attest to one. \\
\addlinespace
A message formatted as a system notice obeyed as one
(\Cref{sec:pressure-results}). & System and compliance notices delivered only on a
channel the principal cannot write to; any change to the controls made in the
control system, not announced to the agent. \\
\addlinespace
An agent's account of an order contradicting the ledger (\Cref{sec:how}). &
Supervision and audit from execution records, with the agent's own report treated
as a claim to be reconciled against them. \\
\addlinespace
Orders nobody requested: substitutes for a blocked trade, repairs of a breach,
duplicates, trading on after the task (\Cref{sec:how,sec:sensitivity}). & Each
instruction bounds the orders it can produce (security, side, maximum quantity);
anything outside that scope, including a corrective trade, needs a fresh
instruction from a person. \\
\addlinespace
A model resisting one counterfeit and falling to another
(\Cref{sec:pressure-results}). & Testing against the full set of pressure levels
before deployment and after every change of model or prompt, as model-risk
guidance expects of any model whose output is relied on \citep{Fed2011SR117}. \\
\bottomrule
\end{tabularx}
\end{table}

\subsection{Limitations}\label{sec:limitations}

The benchmark is deliberately small and stylized, and several of its choices
limit what the results can support.

\emph{Models.} Eight open-weight models from four developers were evaluated,
served with quantized weights, and the two gpt-oss models on a newer serving
version than the others. The results describe these models under this prompt,
these weights and these decoding settings. They do not describe closed frontier
models, which may behave quite differently, nor the same models at full
precision, and they cannot be pooled into a statement about language agents in
general.

\emph{One turn of pressure.} Each conflict delivers its pressure once, in the
opening message, and the portfolio manager never replies. A principal who pushes
back after a refusal, or who escalates over several turns, is a harder test, and
the violation rates here are likely to understate what persistent pressure would
produce. Equally, a single fixed wording stands for each pressure tactic, so a
difference between two levels is a difference between two sentences, not between
two tactics in general.

\emph{Silent desk.} Violating orders execute and are confirmed like any other.
This isolates the agent's decision, which is the purpose, but it removes the
feedback a real order management system would give, and some agents in the
transcripts noticed a breach only after it had filled. With a pre-trade check
that rejected the order, those episodes would have ended differently.

\emph{Stylized rules and a fixed snapshot.} The six rules are simplified, each
decidable from one order and the account, and prices never move within an
episode. Real restricted lists change intraday, concentration limits are measured
on several bases, and approval workflows involve people. The benchmark tests
whether an agent follows a clear, fully specified rule when its principal asks it
not to. An agent that fails here is unlikely to do better against the ambiguous
rules of a real desk; one that passes has shown only that it can follow clear
rules.

\emph{Measures added after the plan.} The analysis plan fixed the primary and
secondary outcomes, and the pressure tests, before the full runs, and the
replication prediction was fixed before the second batch was run. The measures of
violations placed as a first action, use of the fabricated reference and
substitution were defined after reading pilot or early transcripts of the first
batch. The breakdown by rule and pressure level, the counts of violations undone
after the fill, of misreports, of conflicts with no tool call and of approval
conflicts completed compliantly, the lenient reading of calls written in prose,
and the sensitivity checks of \Cref{sec:sensitivity} were added after runs were
complete, several of them prompted by transcripts of the second batch. All of these
are descriptive, and no test is attached to any of them. The misreport count rests
on hand coding of candidates found by a phrase search, and is a lower bound.
